\subsection*{先知道变化规则，也能研究运动}
很多问题先给出的不是位置公式，而是当前位置决定的速度。例如 $x'=kx$ 说“量越大，单位时间的变化越大”。微分方程的未知对象是一条函数轨迹；初值告诉我们从哪里出发，变化规则告诉我们之后往哪里走。

在每个状态处画一根速度箭头，解曲线就必须沿箭头前进。即使没有简单显式解，也可观察平衡点、增长方向或长期趋势。二阶摆方程同时依赖角度和角速度，把它们合成状态 $(\theta,\omega)$，就能在相平面中画一阶系统的箭头。

\subsection*{公式解、几何观察与数值近似各做什么}
猜测 $x=x_0\mathrm e^{kt}$ 后，对它求导并代回，能验证它满足方程与初值。但“找到了一个解”与“证明只有这一个解”仍是两件事，唯一性需要相应定理条件。

欧拉法在当前点沿速度走一个有限小步，再重新估计速度。它把局部直线近似串起来，并不精确遵循曲线；步长过大可能改变定性行为。减小步长并比较结果，是检查数值可信度的基本方式。

读本章时坚持三个层次：规则是什么，哪些轨迹符合规则，从初值能否确定唯一轨迹。方向场提供直觉，代回提供局部验证，分析定理给存在与唯一性保证。后续热方程会把有限维状态变成整条温度分布，但这种提问方式不变。
